%% CPP 2026 Submission — Lean Formalization of the Riemann Direction
%% Compile: pdflatex + bibtex (acmart)
\documentclass[sigplan,10pt,anonymous,review]{acmart}
\settopmatter{printfolios=true,printacmref=false}

\AtBeginDocument{%
  \providecommand\BibTeX{{\normalfont B\scshape ib\TeX}}}

\settopmatter{printacmref=false}
\renewcommand\footnotetextcopyrightpermission[1]{}
\pagestyle{plain}

\usepackage[utf8]{inputenc}
\usepackage[T1]{fontenc}
\usepackage{amsmath}
\ifdefined\XeTeXinterchartokenstate
  % XeLaTeX: unicode-math (loaded by acmart) supplies the amssymb glyphs
\else
  \usepackage{amssymb}
\fi
\usepackage{booktabs}
\usepackage{microtype}

\begin{document}

\title{Formalizing the Riemann Direction in Lean:\\
Projection-Induced Structure Loss, the Circle of Primes,\\
and the Euler Product (C011--C025, all PROVED, no \texttt{sorry})}

\author{Anonymous}
\affiliation{%
  \institution{Independent}
  \country{}}


\date{2026-08-12}

\begin{abstract}
We report a complete formalization, in Lean 4 + mathlib, of a chain of classical
results in the direction of the Riemann hypothesis (claims C011--C025; all PROVED;
novelty: KNOWN; no \texttt{sorry}; full \texttt{lake build} passes). The chain is
intuition-guided: starting from a projection construction of the complex plane
(a two-dimensional rotation algebra projected onto a real axis), through the circle
structure of primes (sums of two squares, lattice points on a circle, uniqueness of
decomposition via Gaussian-integer UFD, splitting into conjugate Gaussian primes),
to the convergence of the Euler product ($\prod_p (1-p^{-s})^{-1} = \sum_n 1/n^s$
for $\Re s > 1$) and the zero-free region ($\Re s \ge 1$). We extract a general
methodological principle---\emph{projection-induced structure loss}---formalized as
the recoverability theorem: lost structure is not recoverable; symmetry directions
are recoverable. The Riemann hypothesis itself is not proved; the critical-line
geometry is an equivalent restatement of the symmetry axis of the functional
equation. The Lean development is fully buildable (mathlib v4.32.2), all theorems
are machine-checked, and the artifact is available.
\end{abstract}

\keywords{Lean, formalization, Riemann zeta, Gaussian integers, Euler product, projection, intuition}

\maketitle

\section{Introduction}

The Riemann zeta function $\zeta(s)$ and its nontrivial zeros are among the most
studied objects in mathematics. The Riemann hypothesis---that all nontrivial zeros
of $\zeta$ lie on the critical line $\Re s = 1/2$---has resisted proof for over 160
years. In parallel, formalization of analytic number theory has advanced: mathlib
now defines $\zeta$, its analytic continuation, the functional equation, and the
statement \texttt{RiemannHypothesis} (unproved).

This paper contributes a \emph{complete, machine-checked} Lean formalization of the
surrounding infrastructure of the Riemann direction: the projection construction of
the complex plane, the circle structure of primes, and the Euler product with its
zero-free region. All results are classical (novelty: KNOWN); the contribution is the
formalization itself and the extracted methodological principle of projection-induced
structure loss.

\paragraph{Contributions.}
\begin{enumerate}
\item A complete Lean formalization (C011--C025, no \texttt{sorry}) of the
infrastructure of the Riemann direction (\S\S3--5).
\item The \emph{recoverability theorem}: projection irrecoverably drops structure
(the imaginary direction) while preserving symmetry directions (the real axis);
formalized as \texttt{proj\_not\_recoverable}, \texttt{proj\_recoverable\_symmetry},
\texttt{projection\_recovery\_theorem}.
\item A buildable artifact (Lean 4, mathlib v4.32.2, full \texttt{lake build}).
\end{enumerate}

\paragraph{Honest boundary.} The Riemann hypothesis is not proved. All results are
restatements of classical mathematics; the critical-line geometry is an algebraic
restatement of the functional-equation symmetry, not a claim about zeros.

\section{The ComplexAxis framework}

\begin{definition}[ComplexAxis]\label{def:axis}
\texttt{ComplexAxis} is the pair $\{\langle a, b\rangle : a,b \in \mathbb{R}\}$
with multiplication $(a_1+b_1J)(a_2+b_2J) = (a_1a_2-b_1b_2) + (a_1b_2+a_2b_1)J$,
$J = \langle 0,1\rangle$; it is isomorphic to the matrix representation of
$\mathbb{C}$.
\end{definition}

\begin{lemma}[$J^2 = -1$, C011]\label{lem:jsq}
$J \cdot J = -1$. In the higher-dimensional structure, $-1$ has a square root.
\end{lemma}

\begin{definition}[projection, lifting]\label{def:proj}
$\texttt{proj}\,\langle a,b\rangle = a$ (drops the rotation component);
$\texttt{lift}\, t = \langle t, 0\rangle$ (real-axis embedding).
\end{definition}

\begin{lemma}[projection drops structure, C011]\label{lem:projdrop}
\texttt{proj} preserves addition but not multiplication: $\texttt{proj}(J\cdot J) = -1
\neq \texttt{proj}(J)\cdot\texttt{proj}(J) = 0$.
\end{lemma}

\begin{lemma}[basepoint drift, C011]\label{lem:drift}
The basepoint is $J$; $\texttt{proj}\, i = 0$ (origin illusion); all purely-imaginary
basepoints project to $0$; basepoint drift is unobservable under projection.
\end{lemma}

\begin{lemma}[real axis is a projection class, C011]\label{lem:realclass}
The real-direction line through any purely-imaginary basepoint projects to $\mathbb{R}$;
the position of the real axis is unobservable under projection.
\end{lemma}

\section{The circle structure of primes}

\begin{theorem}[sums of two squares, C014]\label{thm:twosq}
A prime $p \not\equiv 3 \pmod 4$ is the norm of some point of ComplexAxis:
$\exists a,b,\ a^2+b^2 = p$.
\end{theorem}
\emph{Formalization.} Reuses mathlib's \texttt{Nat.Prime.sq\_add\_sq} (Fermat's
two-square theorem), lifted via \texttt{exact\_mod\_cast}.

\begin{theorem}[unique orbit, C017]\label{thm:unique}
For a prime $p \equiv 1 \pmod 4$, the circle $x^2+y^2=p$ has exactly 8 lattice
points: the representation is unique up to sign and order.
\end{theorem}
\emph{Formalization.} Gaussian-integer UFD: \texttt{gauss\_irreducible\_of\_norm\_prime},
\texttt{gauss\_unit\_of\_norm\_one}, Euclid's lemma. Lean:
\texttt{prime\_sq\_add\_sq\_unique}.

\begin{theorem}[orbit structure, C015--C016]\label{thm:orbit}
The $90^\circ$ rotation ($\times J$) is a 4-cycle preserving norm and lattice; the 8
points are 4 conjugate pairs; lattice points are closed under multiplication; norm is
multiplicative.
\end{theorem}

\begin{theorem}[prime-circle product, C023]\label{thm:product}
The product of the 8 lattice points is $p^4$; $J^2 = -1$ (half-turn), the 4-cycle
closes.
\end{theorem}

\begin{theorem}[splitting, C024]\label{thm:split}
A prime $p \equiv 1 \pmod 4$ splits into a conjugate Gaussian-prime pair
$p = \pi\bar\pi$; the 8 points are the associates of $\pi$ union those of $\bar\pi$.
\end{theorem}

\section{Curling and critical-line geometry}

\begin{theorem}[curling, C015]\label{thm:curl}
The inversion $\texttt{recip}\, z = \bar z/|z|^2$ curls infinity to finiteness: for
every $\varepsilon > 0$ there is $R$ with $|z|>R \Rightarrow |\texttt{recip}\,z|<\varepsilon$;
$\texttt{recip}$ is an involution; $|\texttt{recip}\,z| = 1/|z|$.
\end{theorem}

\begin{theorem}[critical line is a circle, C019]\label{thm:circle}
The critical line ($x=1/2$) is, under inversion, the circle centered at $(1,0)$ with
radius 1. Bidirectional containment with the zero-position set establishes the same
object (C022).
\end{theorem}
\emph{Formalization.} \texttt{critical\_line\_is\_circle},
\texttt{critical\_line\_in\_circle}, \texttt{circle\_recip\_proj},
\texttt{zeroSet\_in\_criticalCircle}, \texttt{criticalCircle\_subset\_zeroSet\_image}.

\begin{theorem}[recoverability, C011]\label{thm:recover}
\emph{Lost structure is not recoverable; symmetry directions are recoverable.}
(i) $i$ and $-i$ project identically ($\texttt{proj}\langle 0,1\rangle =
\texttt{proj}\langle 0,-1\rangle = 0$); (ii) real-axis symmetry is preserved
($\texttt{proj}(\texttt{lift}(-r)) = -\texttt{proj}(\texttt{lift}\,r)$; \texttt{lift}
is injective).
\end{theorem}
\emph{Formalization.} \texttt{proj\_not\_recoverable},
\texttt{proj\_recoverable\_symmetry}, \texttt{projection\_recovery\_theorem}.

\section{Euler product and zero relations}

\begin{theorem}[Euler product, C025]\label{thm:euler}
For $\Re s > 1$: $\prod_p (1-p^{-s})^{-1} = \sum_n 1/n^s$. The identity identifies
with mathlib's \texttt{riemannZeta}, and $\Re s \ge 1$ is zero-free.
\end{theorem}
\emph{Formalization.} \texttt{eulerProduct\_completely\_multiplicative\_tprod},
\texttt{Complex.summable\_one\_div\_nat\_cpow},
\texttt{zeta\_euler\_product}, \texttt{riemannZeta\_euler\_product},
\texttt{riemannZeta\_ne\_zero\_of\_one\_le\_re}.

\section{Methodology: projection-induced structure loss}

The recoverability theorem (Thm.~\ref{thm:recover}) grounds a methodology: by
deliberately constructing a projection that drops the divergent/uncountable structure,
a fast intuition path is obtained at reduced inference cost. In the Riemann direction,
the projection drops the imaginary (analytic-content) direction while preserving the
real-axis symmetry; the retained structure is navigable by the compiled intuition
channel. The author's original framing is the \emph{Worn-Zhe-Yue method}
(the basepoint-construction-induced worn-fold-cross method for mathematical space): ``piercing through, folding, and transcending of mathematical space''
transcending of mathematical space).

The token economy of the formalization ($\approx$700k tokens, 99.2\% context-cache
retransmission, net new content $<1\%$) supports the claim that intuition-guided
formalization lowers inference cost.

\section{Related work}

\subsection{The Riemann direction}
Riemann (1859) introduced the zeta function and the hypothesis; Hardy (1914) proved
infinitely many zeros on the critical line; Levinson (1974) and Conrey (1989) gave
positive proportions. These are external (not formalized here); our formalization
organizes the surrounding geometry.

\subsection{Formalization of analytic number theory in Lean}
mathlib formalizes \texttt{riemannZeta}, the functional equation, and
\texttt{RiemannHypothesis} (unproved). Our work connects the Euler product to the
official \texttt{riemannZeta}.

\subsection{Projection and structure loss}
The recoverability dichotomy (Thm.~\ref{thm:recover}) is the formal version of
projection-induced structure loss, connecting to the fast-path account of the
Unified Framework of Representation, Logic, and Intuition (companion work).

\section{Conclusion}

We report a complete, machine-checked Lean formalization of the infrastructure of the
Riemann direction (C011--C025, no \texttt{sorry}): the projection construction, the
circle structure of primes, the Euler product, and the recoverability dichotomy. The
Riemann hypothesis is not proved; the formalization organizes and verifies the
surrounding geometry.

\paragraph{Future work.} Formalize the partial results on zero distribution (Hardy/
Levinson/Conrey) as mathlib theorems; extend the Worn-Zhe-Yue method to other
projection-induced fast paths.


\bibliographystyle{ACM-Reference-Format}
\bibliography{refs}

\end{document}
